% =============================================================================
% BLOCO 6 — TÉCNICAS QUALITATIVAS E MÉTODOS MISTOS  (17:05–17:50)   [Atividade 6]
% =============================================================================
\secao{Bloco 6}{Técnicas qualitativas e métodos mistos}{17:05 -- 17:50 \,$\cdot$\, Atividade 6}

\begin{frame}{Quantificar o impacto $\times$ compreender o impacto}
  \begin{columns}[T]
    \begin{column}{0.5\textwidth}
      \begin{block}{Abordagem quantitativa}
        \small
        \begin{itemize}
          \item \textbf{Quanto?} Magnitude do \emph{efeito médio};
          \item Foco no que é \emph{geral}, comum;
          \item Lida com viés de seleção; busca o contrafactual;
          \item Generalização, replicação, transparência dos métodos.
        \end{itemize}
      \end{block}
    \end{column}
    \begin{column}{0.5\textwidth}
      \begin{atividadebox}{Abordagem qualitativa}
        \small
        \begin{itemize}
          \item \textbf{Como e por quê?} Qualidade da mudança;
          \item Foco no que é \emph{específico}, no contexto e nos \textbf{significados} que os atores atribuem;
          \item Poucos casos, em profundidade;
          \item Valores, símbolos, percepções.
        \end{itemize}
      \end{atividadebox}
    \end{column}
  \end{columns}
  \vspace{0.2cm}
  \begin{center}\small ``Respondem a diferentes questões'': o quantitativo \emph{identifica}, o qualitativo \emph{explica}.\end{center}
  \fonte{Batista e Domingos (2017), pp.~13--14, a partir de Patton (2002).}
\end{frame}

\begin{frame}{Por que incluir a perspectiva dos atores?}
  \begin{enumerate}
    \item Parte da \ind{perspectiva dos atores} para \textbf{reconstruir a política} (como foi implementada e percebida);
    \item Identifica \ind{impactos diferenciados} sobre diferentes grupos e indivíduos;
    \item Analisa \ind{dimensões não diretamente quantificáveis} (ex.: empoderamento, capital social);
    \item Explica \ind{por que} em algumas situações as políticas atingem os objetivos e em outras não: ajuda a ``abrir a caixa-preta''.
  \end{enumerate}
  \vspace{0.3cm}
  \begin{block}{Uma metáfora do artigo (Patton, 2002)}
    \small Se você quer saber \emph{quanto} pesam as pessoas, use uma balança. Se quer saber \emph{o que o peso significa} para elas, é preciso ouvir suas histórias.
  \end{block}
  \fonte{Batista e Domingos (2017), pp.~14--15.}
\end{frame}

\begin{frame}{Observação e etnografia}
  \begin{columns}[T]
    \begin{column}{0.55\textwidth}
      \small
      \begin{itemize}
        \item \textbf{Imersão} do pesquisador na cultura/grupo; notas de campo, descrições densas, contexto;
        \item \textbf{Participante} (torna-se membro funcional) ou \textbf{não participante} (apenas observa);
        \item O \ind{tempo} é fator central: observações curtas geram interpretações parciais; longos períodos reduzem a \emph{reflexividade} (o observado se comporta ``diferente'' na presença do pesquisador);
        \item Na avaliação: como a política é percebida no cotidiano e \emph{como as pessoas reagem aos seus instrumentos}.
      \end{itemize}
    \end{column}
    \begin{column}{0.45\textwidth}
      \begin{block}{Adato (2008): PTCRs na Nicarágua e na Turquia}
        \footnotesize
        Surveys + etnografia + entrevistas. O survey dizia que o programa era bem \emph{focalizado}, mas a \textbf{comunidade percebia o critério como arbitrário}; mães não davam o suplemento de ferro às crianças por julgarem que \emph{fazia mal}; crianças eram superalimentadas antes da pesagem obrigatória para ``passar no teste''.
      \end{block}
    \end{column}
  \end{columns}
  \fonte{Batista e Domingos (2017), pp.~15--16.}
\end{frame}

\begin{frame}{Entrevistas em profundidade}
  \begin{columns}[T]
    \begin{column}{0.55\textwidth}
      \begin{itemize}
        \item A \emph{palavra} do participante é o foco; perguntas \textbf{abertas e flexíveis}, em oposição aos questionários fechados;
        \item Tipos: \textbf{estruturada}, \textbf{semiestruturada} (tópicos-guia), \textbf{não estruturada} e \textbf{narrativa};
        \item O entrevistador precisa construir \ind{confiança} e preparar um tópico-guia com progressão lógica;
        \item Útil para aprofundar um tema, mapear percepções individuais ou temas \emph{sensíveis}.
      \end{itemize}
    \end{column}
    \begin{column}{0.45\textwidth}
      \begin{block}{Patton (2002): programa de alfabetização de adultos}
        \footnotesize
        O survey registrou que \textbf{77\% estavam ``muito satisfeitos''}. As entrevistas mostraram \emph{por quê}: ler os frascos de remédio das crianças, ler a Bíblia, fazer compras --- ênfase no \textbf{empoderamento} cotidiano, não na colocação no mercado de trabalho (objetivo principal do programa).
      \end{block}
    \end{column}
  \end{columns}
  \fonte{Batista e Domingos (2017), pp.~16--17.}
\end{frame}

\begin{frame}{Grupos focais}
  \begin{columns}[T]
    \begin{column}{0.55\textwidth}
      \begin{itemize}
        \item \textbf{Não} é a soma de entrevistas individuais: o foco está na \ind{discussão e deliberação} entre os participantes;
        \item O moderador \emph{fomenta} a discussão, limita dominantes, incentiva quem fala pouco e mantém o tema;
        \item Bom para explorar o \emph{espectro} de atitudes e opiniões e observar \textbf{consenso e divergência}.
      \end{itemize}
    \end{column}
    \begin{column}{0.45\textwidth}
      \begin{block}{Hunter e Sugiyama (2014): Bolsa Família}
        \footnotesize
        11 grupos focais em 3 municípios do Nordeste (Camaragibe, Jaboatão dos Guararapes e Pau-Brasil) + entrevistas com gestores. Sugerem que o programa deu \textbf{dignidade} às famílias (compra de comida, roupas e material escolar), autonomia frente ao clientelismo e senso de \textbf{empoderamento}.
      \end{block}
    \end{column}
  \end{columns}
  \fonte{Batista e Domingos (2017), pp.~17--18.}
\end{frame}

\begin{frame}{Críticas comuns --- e cuidados}
  \begin{columns}[T]
    \begin{column}{0.5\textwidth}
      \begin{block}{Críticas}
        \small
        \begin{itemize}
          \item \alert{Subjetividade} e especificidade dos resultados;
          \item Ausência de \emph{regularidades} ou generalização;
          \item Menor precisão;
          \item Dificuldade de \emph{replicação}.
        \end{itemize}
      \end{block}
    \end{column}
    \begin{column}{0.5\textwidth}
      \begin{block}{Cuidados e respostas}
        \small
        \begin{itemize}
          \item \textbf{Representatividade} como \ind{variabilidade} (incluir posições diferentes), não como probabilidade estatística;
          \item Selecionar casos usando evidência \emph{quantitativa} (evita viés);
          \item Tempo de campo, treinamento do entrevistador, tópico-guia;
          \item Menos intrusiva: considera a percepção dos próprios atores.
        \end{itemize}
      \end{block}
    \end{column}
  \end{columns}
  \fonte{Batista e Domingos (2017), pp.~14, 18--19.}
\end{frame}

\begin{frame}{Combinando métodos: o que se ganha}
  \begin{columns}[T]
    \begin{column}{0.5\textwidth}
      \begin{enumerate}
        \item \textbf{Validade e credibilidade}: convergência ao medir o mesmo fenômeno (\emph{triangulação});
        \item \textbf{Compreensividade}: capta diferentes dimensões sob perspectivas distintas;
        \item \textbf{Descoberta} de fatores inesperados;
        \item \textbf{Diversidade} de olhares.
      \end{enumerate}
    \end{column}
    \begin{column}{0.5\textwidth}
      \begin{beamercolorbox}[wd=\textwidth,rounded=true,sep=0.3cm]{laranjabox}
        \small
        \textbf{Regra de ouro:} nenhum método é sempre o mais apropriado. O método deve ser escolhido pela \textbf{pergunta} que se quer responder.
        \par\raggedleft\scriptsize Ideia de Worthen, Sanders e Fitzpatrick (1997), citados no artigo
      \end{beamercolorbox}
    \end{column}
  \end{columns}
  \vspace{0.3cm}
  \begin{center}
    \small Greene, Benjamin e Goodyear (2001): quatro contribuições do multimétodo para a avaliação.
  \end{center}
  \fonte{Batista e Domingos (2017), pp.~18--19.}
\end{frame}

\begin{frame}{Coordenados ou integrados?}
  \centering
  \begin{tikzpicture}[scale=0.9,every node/.style={transform shape}]
    % coordenado
    \node[font=\small\bfseries,text=indigo,anchor=west] at (0,3.3) {Coordenado: combinação no final do estudo};
    \node[cxi,minimum width=3.4cm,minimum height=0.8cm,anchor=west] (q1) at (0,2.4) {Método predominante};
    \node[cxl,minimum width=3.4cm,minimum height=0.8cm,anchor=west] (q2) at (6.6,2.4) {Método complementar};
    \draw[setal] (q1) -- node[above=1pt,font=\scriptsize]{ilustra, explica, refina} (q2);
    % integrado
    \node[font=\small\bfseries,text=indigo,anchor=west] at (0,1.2) {Integrado: pontos de interação ao longo de todo o estudo};
    \node[cx,minimum width=2.6cm,minimum height=0.9cm,anchor=west,align=center] (ia) at (0,0.1) {Desenho dos\\instrumentos};
    \node[cx,minimum width=2.6cm,minimum height=0.9cm,anchor=west,align=center] (ib) at (3.4,0.1) {Amostra};
    \node[cx,minimum width=2.6cm,minimum height=0.9cm,anchor=west,align=center] (ic) at (6.8,0.1) {Coleta de\\dados};
    \node[cx,minimum width=2.6cm,minimum height=0.9cm,anchor=west,align=center] (id) at (10.2,0.1) {Interpretação};
    \draw[seta,<->] (ia.east) -- (ib.west);
    \draw[seta,<->] (ib.east) -- (ic.west);
    \draw[seta,<->] (ic.east) -- (id.west);
    \node[font=\scriptsize,text=cinza,anchor=west] at (0,-0.9) {Os métodos quantitativo e qualitativo se alimentam mutuamente em cada etapa.};
  \end{tikzpicture}
  \vspace{0.2cm}
  \small
  \begin{itemize}
    \item Predominantemente quantitativo: o qualitativo \emph{formula instrumentos} e \emph{explica} os resultados;
    \item Predominantemente qualitativo: o quantitativo \emph{seleciona casos} e mostra que eles são representativos dos grupos.
  \end{itemize}
  \fonte{Batista e Domingos (2017), pp.~18--19; Figura 8, adaptada de Bamberger (2012); Caracelli e Greene (1997).}
\end{frame}

\begin{frame}{Exemplos de métodos mistos no artigo}
  \begin{columns}[T]
    \begin{column}{0.33\textwidth}
      \begin{block}{Chile: focalização}
        \scriptsize
        Clert e Wodon (2001): \textbf{survey} nacional mostra grande sucesso na focalização. \textbf{Grupos focais} revelam o que o número escondia: beneficiários não entendem as regras e sentem \emph{injustiça} na distribuição.
      \end{block}
    \end{column}
    \begin{column}{0.33\textwidth}
      \begin{block}{Uruguai: exclusão social}
        \scriptsize
        Baker (2001): \textbf{métodos quantitativos} escolhem as áreas geográficas e os grupos vulneráveis; \textbf{entrevistas e grupos focais} exploram as dimensões e percepções da exclusão. Qualitativo predominante.
      \end{block}
    \end{column}
    \begin{column}{0.33\textwidth}
      \begin{block}{Índia: desenvolvimento comunitário}
        \scriptsize
        Bamberger (2012), abordagem \textbf{integrada}: 200 comunidades sorteadas (tratamento) $\to$ pesquisa exploratória $\to$ survey $\to$ estudo em profundidade de 5 projetos tratados e 5 controles $\to$ \textbf{triangulação}.
      \end{block}
    \end{column}
  \end{columns}
  \fonte{Batista e Domingos (2017), pp.~19--20.}
\end{frame}

\begin{frame}{Atividade 6 \,$\cdot$\, Do número ao porquê (15 min, em grupos)}
  \begin{atividadebox}{Cenário}
    \small
    O DiD do \textbf{Qualifica} mostrou \textbf{+R\$ 300}. Mas: (1) o efeito é \emph{muito menor} entre mulheres com filhos pequenos (R\$~80); (2) 35\% dos inscritos abandonaram o curso; (3) 80\% se dizem ``muito satisfeitos'', mas o gestor acha que ``algo não bate''.
  \end{atividadebox}
  \vspace{0.1cm}
  \small
  \textbf{Em grupos de 3--4 pessoas (salas simultâneas):}
  \begin{enumerate}
    \item \textbf{Parte A}: qual técnica responde a cada pergunta (impacto, motivo das diferenças, rotina do curso, significado de ``satisfeito'')?
    \item \textbf{Parte B}: escrevam \textbf{4 perguntas abertas} para entrevistar mulheres que abandonaram o curso;
    \item \textbf{Parte C}: o desenho de vocês é \emph{coordenado} ou \emph{integrado}? Que risco metodológico vocês enxergam?
  \end{enumerate}
  \vfill{\scriptsize Roteiro completo: \cod{atividades/atividade\_06\_qualitativo\_metodos\_mistos.md}. Plenária: 10 min.}
\end{frame}
